\documentclass[10pt,a4paper]{amsart}
\usepackage[margin=19mm]{geometry}
\usepackage{amsmath,amssymb,mathtools}
\usepackage[scr=rsfs,cal=euler]{mathalfa}
\usepackage{xfrac}
\usepackage[unicode,hidelinks,bookmarks=false]{hyperref}
\newcommand{\ie}{i.e.\ }
\newcommand{\cf}{cf.\ }
\newcommand{\resp}{respectively\ }
\newcommand{\Subsection}{\subsection}
\providecommand{\nobreakdash}{\nobreak\hbox{-}}
\setlength{\emergencystretch}{2em}
\multlinegap0pt
\usepackage{amssymb}
\newtheorem{lemma}{Lemma}
\newtheorem{proposition}{Proposition}
\newcommand{\R}{\mathbb R}
\title{Harnack Inequality for Parabolic Fractional $p$-Laplace Equations}
\author{Harsh Prasad}
\date{Source: arXiv:2609.15694v1 (CC BY 4.0)}
\begin{document}
\maketitle

\noindent\emph{Source and licence.} Harnack Inequality for Parabolic Fractional $p$-Laplace Equations. Version arXiv:2609.15694v1 (CC BY 4.0), distributed by arXiv under the Creative Commons Attribution 4.0 International licence, http://creativecommons.org/licenses/by/4.0/, as recorded in the arXiv licence metadata for this exact version and shown on the abstract page https://arxiv.org/abs/2609.15694v1. Full licence notice: \texttt{licenses/arxiv-2609.15694-CC-BY-4.0.txt}.\par\smallskip

\noindent\textit{Editorial scope.} Counted unit: the article's proposition prop:moment (source0.tex lines 896--919). The article's complete original proof of it is delivered with it (source0.tex lines 921--1096, one \texttt{proof} environment of 176 lines). No mathematics is written, repaired or re-derived: the statement and the proof are the article's own lines, in the article's own notation, and no retained line is substituted or re-flowed.  Retained uncounted as the source-local context the proof needs: lines 349--355; lines 756--800.  Nothing was omitted from any retained span, so the package carries no omission record.  No retained line contains a citation, so the wrapper carries no bibliography and no bibliography entry.  No retained line contains a figure environment, an image-inclusion command or drawing code, so no image is delivered and the package declares no asset dependency.  Editorial formatting only, in addition: an AMS \texttt{amsart} preamble that restates the article's own declarations and macros used by the retained lines, namely the environments \texttt{lemma}, \texttt{proposition} on separate counters, together with the standard packages those lines need.  The retained environments print in this document as Lemma 1, Proposition 1 in order of appearance, whereas the article prints the same items with the numbering of its own full document.  The complete native source of the article as distributed in the arXiv e-print for this exact version is bundled unchanged as \texttt{sources/210441/source0.tex}.
\begin{equation}\label{eq:kernel}
K(x,y,t)=K(y,x,t),
\qquad
\lambda |x-y|^{-N}
\le K(x,y,t)
\le \Lambda |x-y|^{-N}.
\end{equation}

\begin{lemma}\label{lem:continuity}
There exist constants

\[
    \delta_0>0,
    \qquad
    \rho_0\in(0,1),
    \qquad
    c_0\in(0,1),
\]

depending only on the data, with the following property. Let $v$ be a locally
bounded weak solution in

\[
    B_4\times(-4^{sp},0]
\]

such that

\[
    0\le v\le2
    \quad\text{in }B_4\times(-4^{sp},0],
    \qquad
    v(0,0)=1.
\]

If
\begin{equation}\label{eq:normalized-tail-small}
\int_{-4^{sp}}^0
\int_{\R^n\setminus B_4}
\frac{|v(y,t)|^{p-1}}{|y|^{n+sp}}\,dy\,dt
\le\delta_0,
\end{equation}
then
\begin{equation}\label{eq:normalized-positive}
v\ge c_0
\end{equation}
in a fixed backward cylinder

\[
    B_{\rho_0}\times(-\rho_0^{sp},0].
\]

\end{lemma}

\begin{proposition}\label{prop:moment}
There exists $\varepsilon_0>0$, depending only on the data, such that the
following holds. Suppose $u$ is a locally bounded weak solution in
$B_2\times(-1,1)$ and assume that

\[
    u\ge0
    \quad\text{a.e. in }B_2\times(-1,0].
\]

Suppose moreover that

\[
    u(0,0)=1.
\]

Then
\begin{equation}\label{eq:absolute-moment}
\int_{-1}^0
\int_{\R^n}
|u(y,t)|^{p-1}\omega(y)\,dy\,dt
\ge\varepsilon_0.
\end{equation}
\end{proposition}

\begin{proof}
We set
\[
\mathcal M
:=
\int_{-1}^0\int_{\R^n}
|u(y,t)|^{p-1}\omega(y)\,dy\,dt.
\]
There is nothing to prove if $\mathcal M=\infty$, so there is no loss of generality in assuming that 
it is finite. We first select a point and an intrinsic cylinder on which the value
of $u$ can at most double.

Indeed, let $c\in(0,1/8)$ be a constant, depending only on the data, sufficiently
small as specified below. We run a stopping time argument. We start with
\[
z_0=(x_0,t_0)=(0,0),
\qquad
M_0=u(z_0)=1,
\qquad
r_0=cM_0^{-1/N}.
\]
Given $z_j=(x_j,t_j)$ and $M_j=u(z_j)$, put
\begin{equation}\label{eq:doubling-cylinder}
\mathcal Q_j
:=
B_{4r_j}(x_j)
\times
\bigl(t_j-M_j^{2-p}(4r_j)^{sp},t_j\bigr],
\qquad
r_j:=cM_j^{-1/N}.
\end{equation}
If $u\le2M_j$ in $\mathcal Q_j$, the construction stops. Otherwise, by
continuity, we may choose $z_{j+1}\in\mathcal Q_j$ such that
$M_{j+1}:=u(z_{j+1})>2M_j$.

Let
\[
\alpha:=p-2+\frac{sp}{N}>0.
\]
As long as the construction continues, $M_j>2^j$, and hence
\begin{align*}
|x_j|
&\le
4c\sum_{i=0}^{j-1}M_i^{-1/N}
\le
4c\sum_{i=0}^\infty2^{-i/N},
\\
-t_j
&\le
4^{sp}c^{sp}
\sum_{i=0}^{j-1}M_i^{-\alpha}
\le
4^{sp}c^{sp}
\sum_{i=0}^\infty2^{-i\alpha}.
\end{align*}
We fix $c$ so small that these bounds, including one further cylinder as in
\eqref{eq:doubling-cylinder}, keep every $\mathcal Q_j$ inside
$B_{1/2}\times(-1/2,0]$. The iteration must stop after finitely many steps
because otherwise $M_j>2^j$ at points in a fixed compact subset of
$B_2\times(-1,1)$, contradicting the local boundedness of $u$.

Let $z_*=(x_*,t_*)$, $M=u(z_*)$, and $r=cM^{-1/N}$ be the quantities at the
stopping index. Thus $M\ge1$ and
\begin{equation}\label{eq:doubling-bound}
0\le u\le2M
\quad\text{in }
B_{4r}(x_*)
\times
\bigl(t_*-M^{2-p}(4r)^{sp},t_*\bigr].
\end{equation}
Define the intrinsic rescaling
\begin{equation}\label{eq:moment-rescaling}
v(x,t)
:=
\frac1M
u\bigl(x_*+rx,t_*+M^{2-p}r^{sp}t\bigr).
\end{equation}
The rescaled kernel
\[
K_*(x,y,t)
:=
r^N
K\bigl(x_*+rx,x_*+ry,t_*+M^{2-p}r^{sp}t\bigr)
\]
satisfies \eqref{eq:kernel} with the same constants. Consequently, $v$ is a
locally bounded weak solution in $B_4\times(-4^{sp},0]$, and
\begin{equation}\label{eq:rescaled-local-bounds}
0\le v\le2,
\qquad
v(0,0)=1.
\end{equation}

We compute
\begin{equation}\label{eq:rescaled-tail-identity}
\begin{split}
&\int_{-4^{sp}}^0
\int_{\R^n\setminus B_4}
\frac{|v(y,t)|^{p-1}}{|y|^N}\,dy\,dt
\\
&\qquad=
\frac1M
\int_{t_*-M^{2-p}(4r)^{sp}}^{t_*}
\int_{\R^n\setminus B_{4r}(x_*)}
\frac{|u(z,\tau)|^{p-1}}{|z-x_*|^N}\,dz\,d\tau.
\end{split}
\end{equation}
For $|z-x_*|>4r$, we claim that
\begin{equation}\label{eq:kernel-weight-comparison}
\frac1{M|z-x_*|^N}
\le
C_c\omega(z),
\end{equation}
where $C_c$ depends only on $c$ and the data. Indeed, if $|z|\le2$, then
\[
\frac1{M|z-x_*|^N}
\le
\frac1{M(4r)^N}
=
(4c)^{-N},
\]
whereas $\omega(z)\ge3^{-N}$. If $|z|>2$, then $|x_*|<1$ gives
$|z-x_*|\ge|z|/2$, and the claim follows from $M\ge1$ and
$|z|^{-N}\le C\omega(z)$. Since the time interval in
\eqref{eq:rescaled-tail-identity} lies in $[-1,0]$, we conclude that
\begin{equation}\label{eq:rescaled-tail-moment-bound}
\int_{-4^{sp}}^0
\int_{\R^n\setminus B_4}
\frac{|v(y,t)|^{p-1}}{|y|^N}\,dy\,dt
\le
C_c\mathcal M.
\end{equation}

Let $\delta_0$, $\rho_0$, and $c_0$ be the constants from Lemma
\ref{lem:continuity}. If $\mathcal M\le\delta_0/C_c$, then we get
that
\[
v\ge c_0
\quad\text{in }
B_{\rho_0}\times(-\rho_0^{sp},0].
\]
Returning to the original variables, we have
\begin{equation}\label{eq:physical-positive-cylinder}
u\ge c_0M
\end{equation}
in
\[
B_{\rho_0r}(x_*)
\times
\bigl(t_*-M^{2-p}(\rho_0r)^{sp},t_*\bigr].
\]
This cylinder is contained in $B_1\times[-1,0]$, where
$\omega\ge2^{-N}$. Therefore
\begin{align*}
\mathcal M
&\ge
2^{-N}c_0^{p-1}M^{p-1}
|B_{\rho_0r}|
M^{2-p}(\rho_0r)^{sp}
\\
&=
2^{-N}|B_1|c_0^{p-1}\rho_0^N
Mr^N
\\
&=
2^{-N}|B_1|c_0^{p-1}\rho_0^Nc^N
=:
\varepsilon_1>0.
\end{align*}
If instead $\mathcal M>\delta_0/C_c$, we already have a universal lower
bound. Thus \eqref{eq:absolute-moment} follows with
\[
\varepsilon_0
:=
\min\left\{\frac{\delta_0}{C_c},\varepsilon_1\right\}.
\]
\end{proof}

\end{document}
